\uuid{inO0}
\titre{Espaces supplémentaires/Base/Dimension}
\auteur{hanani ; mzari}
\organisation{exo7}
\licence{CC BY-NC-SA 4.0}
\datecreate{2020-01-20}
\niveau{L1}
\module{Algèbre}
\chapitre{Espace vectoriel}
\sousChapitre{Somme directe}
\difficulte{1}
\theme{Espaces vectoriels supplémentaires, Espaces supplémentaires, Base, Dimension}
\exo7qcmid{lille-2.espaces-vectoriels.053}

\contenu{
\begin{question}
\qclassification{106.03}
\qtags{motcle=Espaces supplémentaires/Base/Dimension}
On considère les deux sous-espaces vectoriels de $\Rr^4$ :
$$E= \{(x,y,z,t) \in \Rr^4 \, ; \; x-y=y-z=t=0\}\; \mbox{ et }\; F= \{(x,y,z,t) \in \Rr^4 \, ; \; z=x+y \}.$$
Quelles sont les assertions vraies ?
\begin{answers}
\good{$\dim E= 1$}.
\bad{$\dim F = 2$}.
\bad{$\dim E\cap F = 1$}.
\good{$E$ et $F$ sont supplémentaires dans $\Rr^4$}.
\end{answers}
\begin{explanations} Une base de $E$ est $\{u\}$, où $u= (1,1,1,0)$, donc $\dim E=1$. Une base de $F$ est $\{v_1,v_2,v_3\}$, où $v_1= (1,0,1,0), v_2= (0,1,1,0)$ et $v_3= (0,0,0,1)$, donc $\dim F=3$. On vérifie que $E\cap F =\{(0,0,0,0)\}$. Donc, d'après le théorème de la dimension d'une somme, $\dim (E+F)=4=\dim \Rr^4$ et comme, en plus, $E+F$ est un sous-espace de $\Rr^4$, $E+F=\Rr^4$. Par conséquent, $E$ et $F$ sont supplémentaires dans $\Rr^4$.
\end{explanations}
\end{question}
}
